% ---------------------------------------------------------------------
% Section 3
% ---------------------------------------------------------------------

\begin{cnread}
本节回答一个自然的问题：切触空间里那两个投影 $\pi$（忘掉赋值 $a$）与
$Q$（忘掉 Reeb 荷 $u$）之间的「勒让德对偶」，究竟与力学里联系
Lagrange 量和 Hamilton 量的\textbf{勒让德变换}是不是同一件事？答案是：
在标准切触流形上，它就是同一个操作的两个 Legendre 叶层。
由此得到本文的第一个真正的新命题：\textbf{双重撕裂定理}——固定荷看
「赋值撕裂 $\Delta a$」（Paper~0 定理 8.9），固定状态看「Reeb 纤维漂移
$\Delta A$」（Paper~0 定理 9.9），两者是同一个切触曲率的两个共轭读法。
本节末尾用核字 $h(p,q)$ 做两条路径的现场验算，结果一字不差是
$p(1-e^{-q})$。
\end{cnread}

\subsection{Why the mechanical Legendre transform is a contact statement}

In analytical mechanics with one degree of freedom, let $q$ be the coordinate,
$p$ the momentum, and $S$ the action, so that $(q,p,S)$ are coordinates on the
first jet bundle $J^1(\R)\cong\R^3$.  The canonical contact form is
\begin{equation}\label{eq:mech-contact}
  \alpha_{\mathrm{mech}}=\mathrm dS-p\,\mathrm dq=0 .
\end{equation}
Using $\mathrm d(pq)=p\,\mathrm dq+q\,\mathrm dp$ and putting
$\tilde S:=pq-S$, the same form reads
\begin{equation}\label{eq:mech-contact-dual}
  \alpha_{\mathrm{mech}}=-\mathrm d\tilde S+q\,\mathrm dp=0 .
\end{equation}
Equations \eqref{eq:mech-contact} and \eqref{eq:mech-contact-dual} are the same
one-form written through two different choices of potential, and the two
choices are related by the Legendre transformation.  Geometrically the content
is that one contact three-manifold carries two Legendrian fibrations: the
``coordinate'' projection $(q,p,S)\mapsto(q,S)$ in which $p\,\mathrm dq$ is read
as the accumulated quantity, and the ``momentum'' projection in which
$q\,\mathrm dp$ is.  Nothing about the contact structure changes; what changes is
which variable is declared to be a function of the others.

\begin{remark}[What the Legendre transformation does not do]
The passage from \eqref{eq:mech-contact} to \eqref{eq:mech-contact-dual} uses
only the product rule.  It is therefore available wherever the product exists;
the analytic content of the classical Legendre transformation (convexity,
invertibility of $p=\partial S/\partial q$) is not used and is not needed here.
\end{remark}

\subsection{The arithmetic contact form and its two potentials}

In the natural charge coordinates $A=\mu u$ and $M=\lambda v$ of Paper~0,
Appendix~D.4, the arithmetic contact form is
\begin{equation}\label{eq:arith-contact}
  \alpha=\mathrm da-\mathrm dA-a\,\mathrm dM,
\end{equation}
and the global affine change of coordinates
\begin{equation}\label{eq:darboux}
  x_{\mathrm D}=M,
  \qquad
  y_{\mathrm D}=a,
  \qquad
  z_{\mathrm D}=a-A
\end{equation}
puts it in standard Darboux form $\alpha=\mathrm dz_{\mathrm D}-y_{\mathrm D}
\,\mathrm dx_{\mathrm D}$ (Appendix~D.4).  Comparing \eqref{eq:darboux} with
\eqref{eq:mech-contact} identifies the correspondence
\begin{equation}\label{eq:legendre-dictionary}
  S\;\longleftrightarrow\; a-A,
  \qquad
  p\;\longleftrightarrow\;a,
  \qquad
  q\;\longleftrightarrow\;M .
\end{equation}
The current assignment $a$ plays the role of the momentum, and the logarithmic
scale $M$ plays the role of the coordinate.

\begin{lemma}[Legendre-dual form of $\alpha$]
\label{lem:legendre-linear}
With $\tilde S:=aM-a+A$ one has the identity
\begin{equation}\label{eq:legendre-linear}
  \alpha=-\mathrm d\tilde S+M\,\mathrm da
  \qquad\text{on }\Cont .
\end{equation}
\end{lemma}

\begin{proof}
Expanding $-\mathrm d(aM-a+A)+M\,\mathrm da$ gives
$-(a\,\mathrm dM+M\,\mathrm da-\mathrm da+\mathrm dA)+M\,\mathrm da
=\mathrm da-\mathrm dA-a\,\mathrm dM$.
\end{proof}

\begin{remark}[A correction to the working note]
\label{rem:legendre-correction}
The working note from which this section is drawn displayed the same step with
the coefficient $1+M$ on $\mathrm da$, that is
$\alpha=-\mathrm d(A+aM-a)+(1+M)\mathrm da$.  Expanding shows that this has an
extra $\mathrm da$ and is not $\alpha$.  The correct dual potential is
$A+aM-a$ with momentum coefficient $M$, as in \eqref{eq:legendre-linear}.  The
error does not propagate: every later identity in the note uses the
exponential form \eqref{eq:legendre-exponential} below, which is correct as
written.  Check~V5.1--V5.2 of Appendix~\ref{app:register} records both.
\end{remark}

The linear Legendre transformation \eqref{eq:legendre-linear} is not the most
intrinsic one available here, because the multiplicative motion of \AEG{} is the
exponential $e^{M}$ rather than the addition of $M$.  Rescaling the contact
form by the nowhere-vanishing factor $e^{-M}$ --- which leaves the horizontal
distribution $\ker\alpha$ untouched --- gives a second, multiplicative dual.

\begin{lemma}[Exponential Legendre identity]
\label{lem:legendre-exponential}
On $\Cont$,
\begin{equation}\label{eq:legendre-exponential}
  e^{-M}\alpha=\mathrm d\bigl(e^{-M}a\bigr)-e^{-M}\mathrm dA .
\end{equation}
\end{lemma}

\begin{proof}
$e^{-M}\alpha=e^{-M}(\mathrm da-a\,\mathrm dM)-e^{-M}\mathrm dA$ and
$\mathrm d(e^{-M}a)=e^{-M}\mathrm da-a e^{-M}\mathrm dM$.
\end{proof}

\begin{remark}[Three names for one quantity]
\label{rem:hat-a}
The quantity $\hat a:=e^{-M}a$ occurring in \eqref{eq:legendre-exponential} has
three readings.  It is the source-normalized assignment of Paper~0,
Proposition~8.2; it is minus the horizontal coordinate of the basic hyperbolic
model, because under the Reeb quotient
$Q(u,v,a)=(-ae^{-\lambda v},e^{-\lambda v})$ one has $\hat a=-x$; and it is, by
Lemma~\ref{lem:legendre-exponential}, the potential whose differential measures
the additive charge in the target frame.  Paper~0, Theorem~9.9 is exactly the
statement that these three readings agree.
\end{remark}

\subsection{The dual tearing theorem}

The two readings of Lemmas~\ref{lem:legendre-linear}
and~\ref{lem:legendre-exponential} become two comparison experiments.  The
statement below is Paper~0, Theorem~9.9; it is reproduced here in the notation
of this note because every later section uses it.

\begin{theorem}[Dual tearing and Reeb drift]
\label{thm:dual-tearing}
Let $\tilde\gamma_1,\tilde\gamma_2$ be piecewise smooth horizontal lifts in
$\Cont$ with a common initial state.
\begin{enumerate}[label=(\arabic*)]
  \item \emph{(Fixed charges: value tearing.)}  Suppose the projected charge
  chains $C_i=\pi\circ\tilde\gamma_i$ share their terminal point $(A_*,M_*)$.
  Then the terminal assignments differ by
  \begin{equation}\label{eq:value-tearing}
    a(\tilde\gamma_1(1))-a(\tilde\gamma_2(1))
    =e^{M_*}\oint_{C_1-C_2}e^{-M}\mathrm dA
    =\iint_{\Sigma_{12}}e^{M_*-M}\mathrm dA\wedge\mathrm dM
    =\tau(\gamma_1,\gamma_2),
  \end{equation}
  for any oriented two-chain $\Sigma_{12}$ in $\ACS$ with
  $\pt\Sigma_{12}=C_1-C_2$.  This is the weighted torsion--Stokes formula of
  Paper~0, Theorem~8.9.
  \item \emph{(Fixed state: Reeb drift.)}  Suppose instead that the projected
  paths $\tilde\gamma_i=Q\circ\tilde\gamma_i$ terminate at the same point
  $p\in\HH^2$, equivalently at the same pair $(M_*,a_*)$.  Then the terminals
  lie on the same Reeb fiber and
  \begin{equation}\label{eq:reeb-drift}
    \Delta A:=A(\tilde\gamma_1(1))-A(\tilde\gamma_2(1))
    =-\oint_{\tilde\gamma_1-\tilde\gamma_2}a\,\mathrm dM
    =-\iint_{\Omega_{12}}\mathrm da\wedge\mathrm dM
    =-\mu\lambda\,\mathrm{Area}_{g_{\mu,\lambda}}(\Omega_{12}),
  \end{equation}
  for any oriented two-chain $\Omega_{12}\subset\HH^2$ with
  $\pt\Omega_{12}=\tilde\gamma_1-\tilde\gamma_2$.
\end{enumerate}
\end{theorem}

\begin{proof}
(1) Integrate \eqref{eq:legendre-exponential} along each horizontal lift; the
initial terms cancel on subtraction, and $e^{M_*}$ multiplying the result gives
the first equality.  The second is Stokes' theorem in $\ACS$ with the target
weight $e^{M_*-M}$, as in Paper~0, Theorem~8.9.

(2) Integrate \eqref{eq:arith-contact} along each lift to get
$A(\tilde\gamma_i(1))-A_0=(a(\tilde\gamma_i(1))-a_0)-\int_{\tilde\gamma_i}
a\,\mathrm dM$.  Under common terminal state the boundary terms cancel on
subtraction, giving the line-integral form.  Stokes' theorem in $\HH^2$ and
$d\alpha=-\mu\lambda\,Q^*\mathrm d\vol_{g_{\mu,\lambda}}$ of Paper~0,
Proposition~9.8(5) give the area form.
\end{proof}

\begin{remark}[Why the two projections must be conjugate]
\label{rem:conjugacy}
$\pi$ forgets the assignment fiber $\pt_a$ and $Q$ forgets the Reeb fiber
$\pt_u$.  These are different directions of one three-manifold, so a comparison
that closes in one projection necessarily fails to close in the other: a
charge-compatible pair stays open on $\HH^2$, and a state-compatible pair closes
on $\HH^2$ and opens along the Reeb fiber.  That complementarity --- not an
analogy --- is the content of the word ``duality'' here.
\end{remark}

\subsection{A worked pair: the kernel word}

The smallest nontrivial instance is the element of $\ker\nu$ exhibited in
Paper~0, Theorem~8.6,
\begin{equation}\label{eq:kernel-word}
  h(p,q)=\cM_q\,\cA_p\,\cM_{-q}\,\cA_{-e^{-q}p},
\end{equation}
whose two halves are the prefixes $\gamma_1=(\cM_q,\cA_p)$ and
$\gamma_2=(\cA_{e^{-q}p},\cM_q)$.  Both carry $(M,a)=(0,0)$ to the same terminal
state with $M_*=q$ and $a_*=p$.

\begin{lemma}[Reeb drift of the kernel word]
\label{lem:kernel-word}
For $\gamma_1,\gamma_2$ as above,
\[
  \Delta A=A(\gamma_1)-A(\gamma_2)=p-e^{-q}p=p(1-e^{-q}),
\]
and the two sides of \eqref{eq:reeb-drift} both evaluate to $p(1-e^{-q})$.  The
weighted torsion of \eqref{eq:value-tearing} for the same pair vanishes, as it
must, since the terminal assignments agree.
\end{lemma}

\begin{proof}
The bare additive charges are $A(\gamma_1)=p$ and $A(\gamma_2)=e^{-q}p$.  On the
state side, $\gamma_1$ is the lower broken line $(0,0)\to(q,0)\to(q,p)$ with
$a\equiv0$ along its horizontal leg, so $\int_{\gamma_1}a\,\mathrm dM=0$, while
$\gamma_2$ follows $a(M)=pe^{M-q}$ and
$\int_{\gamma_2}a\,\mathrm dM=\int_0^q pe^{M-q}\mathrm dM=p(1-e^{-q})$.
Hence $-\bigl(\int_{\gamma_1}-\int_{\gamma_2}\bigr)a\,\mathrm dM=p(1-e^{-q})$,
in agreement with the charge difference.  The weighted area route gives
$\int_0^q\!\!\int_0^p e^{-M}\mathrm dA\,\mathrm dM=p(1-e^{-q})$ as well.  The
verification is Check~V6.1--V6.3.
\end{proof}

This is the identity that Paper~0, Appendix~D.2 records as the archetype of
``zero total charge is weaker than operator neutrality'': the loop
$\ell=(\cA_p,\cM_q,\cA_{-p},\cM_{-q})$ has total charge $(0,0)$ yet
$\tau(\ell,\varnothing)=p(1-e^{-q})$.  In Section~\ref{sec:coulomb} this lemma is
read as the binding energy of a charge dipole, and in
Section~\ref{sec:radiation} it is read as the phase carried by a virtual pair.

\subsection{The duality table}

\begin{center}\small
\renewcommand{\arraystretch}{1.3}
\begin{tabular}{@{}p{0.22\linewidth}p{0.36\linewidth}p{0.36\linewidth}@{}}
\toprule
 & \textbf{reading A: charge projection $\pi$} & \textbf{reading B: Reeb quotient
 $Q$} \\
\midrule
map & $\Cont\to\ACS_{(A,M)}$, forgets $a$
& $\Cont\to\E_0$, forgets $u=A/\mu$ \\
form used & $e^{-M}\alpha=\mathrm d(e^{-M}a)-e^{-M}\mathrm dA$
& $\alpha=-\mathrm dA+(\mathrm da-a\,\mathrm dM)$ \\
closure hypothesis & same terminal charges $(A_*,M_*)$
& same terminal state $p=(M_*,a_*)$ \\
vertical defect & value tearing $\Delta a$
& Reeb drift $\Delta A=\mu\Delta u$ \\
value & $\Delta a=\iint_\Sigma e^{M_*-M}\mathrm dA\wedge\mathrm dM$
& $\Delta A=-\iint_\Omega\mathrm da\wedge\mathrm dM
=-\mu\lambda\,\mathrm{Area}_g(\Omega)$ \\
\AEG{} source & Theorem~8.9, Prop.~8.2
& Theorem~8.6, Prop.~9.8, Theorem~9.9 \\
\bottomrule
\end{tabular}
\end{center}

\subsection{The open thread: from the affine sector to the projective group}

Both readings above live in the affine sector: the horizontal lifts are lifts
inside $B_\infty$, and the Riccati structure of the assignment along a flow is
hidden.  If one allows the second-slot division $c/z$, i.e. the inverting
element $J(z)=-1/z$, then the same one-form acquires a quadratic term and the
assignment obeys a Riccati rather than a linear equation.  Section~\ref{sec:radiation}
carries out exactly that step, and shows that the additional datum which then
appears --- the kernel coordinate $b$ of the bivaluation --- is what makes a
closed history carry a non-telescoping holonomy.
